\documentclass[11pt,a4paper]{article}
\usepackage[margin=24mm]{geometry}
\usepackage{fontspec}
\setmonofont[Scale=.85,Path=paper/v2/fonts/,Extension=.ttf,UprightFont=*,BoldFont=*-Bold,ItalicFont=*-Oblique,BoldItalicFont=*-BoldOblique]{DejaVuSansMono}
\usepackage{amsmath,amssymb,amsthm}
\usepackage{array,booktabs,tabularx,enumitem}
\usepackage[table]{xcolor}
\definecolor{ink}{HTML}{1E293B}
\definecolor{slate}{HTML}{475569}
\definecolor{rulegray}{HTML}{CBD5E1}
\definecolor{tint}{HTML}{F8FAFC}
\definecolor{accent}{HTML}{174B69}
\usepackage[unicode,colorlinks=true,linkcolor=accent,urlcolor=accent,citecolor=accent,bookmarksnumbered=true]{hyperref}
\hypersetup{pdfauthor={Aleksandr Strand}}
\setlength{\emergencystretch}{2.5em}
\setcounter{tocdepth}{1}
\setlist{itemsep=1pt,topsep=3pt}
\newcommand{\src}[2]{\href{https://github.com/Dnarts/rule110-dependency-fronts/blob/v2.2.1-monograph-multilingual/#1}{#2}}

\setmainfont[Path=translations/fonts/,Extension=.otf,UprightFont=*-Regular,BoldFont=*-Regular,BoldFeatures={FakeBold=.8},ItalicFont=*-Regular,ItalicFeatures={FakeSlant=.13},BoldItalicFont=*-Regular,BoldItalicFeatures={FakeBold=.8,FakeSlant=.13}]{ObserverSerifSC}
\linespread{1.10}
\XeTeXlinebreaklocale "zh"\relax
\XeTeXlinebreakskip=0pt plus .1pt\relax
\renewcommand{\contentsname}{目录}
\renewcommand{\refname}{参考文献}
\renewcommand{\proofname}{证明}
\renewcommand{\tablename}{表}
\hypersetup{pdflang={zh-CN},pdftitle={110 OBSERVER：系统自身过去的可观测性},pdfsubject={精确解码、规则 110 与 OLMo 实验}}
\newtheorem{theorem}{定理}[section]
\newtheorem{lemma}[theorem]{引理}
\newtheorem{proposition}[theorem]{命题}
\newtheorem{conjecture}[theorem]{猜想}
\theoremstyle{definition}
\newtheorem{definition}[theorem]{定义}
\newtheorem{remark}[theorem]{注}
\newcommand{\edition}{zh}
\begin{document}
\sloppy
\begin{center}
{\small\bfseries\color{accent}110 OBSERVER}\\[7pt]
{\LARGE\bfseries 系统自身过去的可观测性}\\[5pt]
{\large 消息代价、因果见证与推断的边界}\\[10pt]
Aleksandr Strand \hspace{1em} $\cdot$ \hspace{1em} 2026 年 11 月 20 日
\end{center}
\noindent\textbf{摘要}\quad 系统为作出下一次决策，至少需要多少关于过去的信息？怎样检验它是否使用了这些信息？本文给出精确解码的判据，并确定所需的消息代价。在规则~110 的主格点系统中，两个构型合并后，观察者将保存在外部单比特寄存器中的传感器读数送回格点。系统在下一次更新时恢复过去的比特。在冻结的 OLMo 中，发送方在新串出现之前，将先前的选择与反馈编码为一个规则词元。选择与反馈的四种组合产生四条消息，分别在不同的新串上复用。当首尾比特不同时，接收方用传来的词元答对开发集（dev）中的 64/64 个条件，改为恒定词元 A 时答对 32/64 个。在整个开发集上，答对 124/128 个条件，完整通过 6/8 个区组。对信道的干预检验了规则的传递。要从中对主观体验作出推断，仍需要额外前提。

\begin{quote}\small\raggedright
\textbf{本版验证范围}\quad 无限格点上规则~110 的精确几何定理（\texttt{Rails.essential\_iff}）已在 Lean~4.35.0 中证明。在编写本专著时，还使用两个规范导入模块单独检查了 \texttt{Rule110Observer.lean}。日志记录了全部 11/11 份公理报告（附录~\ref{app:lean}）。OLMo 两步方案在八个中间串上取得 \texttt{PASS\_DEV}。留出面板 9--20、21--32、33--44 和激活信道 G3 的状态均为 \texttt{NOT\_RUN}。桥接公设的状态为 \texttt{DECLARED\_NOT\_DERIVED}。
\end{quote}
\tableofcontents
\clearpage
\input{translations/\edition/01-introduction}
\input{translations/\edition/02-theory}
\input{translations/\edition/03-geometry}
\input{translations/\edition/04-witnesses}
\input{translations/\edition/05-olmo}
\input{translations/\edition/06-boundary}
\input{translations/\edition/07-discussion}
\appendix
\input{translations/\edition/app-a-lean}
\input{translations/\edition/app-b-gates}
\input{translations/\edition/app-c-registers}
\clearpage
{\small\begin{thebibliography}{9}
\bibitem{cook2004}
M.~Cook, ``Universality in elementary cellular automata,'' \emph{Complex Systems}, vol.~15, no.~1, pp.~1--40, 2004.

\bibitem{martinez2006}
G.~J.~Mart\'inez, H.~V.~McIntosh, and J.~C.~Seck-Tuoh-Mora, ``Gliders in Rule 110,'' \emph{International Journal of Unconventional Computing}, vol.~2, no.~1, pp.~1--49, 2006.

\bibitem{shalizi2001}
C.~R.~Shalizi and J.~P.~Crutchfield, ``Computational mechanics: Pattern and prediction, structure and simplicity,'' \emph{Journal of Statistical Physics}, vol.~104, pp.~817--879, 2001.

\bibitem{wolfram1986}
S.~Wolfram, ``Tables of cellular automaton properties,'' in \emph{Theory and Applications of Cellular Automata}, World Scientific, pp.~485--557, 1986.

\bibitem{wolpert2008}
D.~H.~Wolpert, ``Physical limits of inference,'' \emph{Physica D: Nonlinear Phenomena}, vol.~237, no.~9, pp.~1257--1281, 2008.
\bibitem{witsenhausen}
H.~S. Witsenhausen. The zero-error side information problem and chromatic numbers.
\emph{IEEE Transactions on Information Theory}, 22(5), 592--593, 1976.
\href{https://doi.org/10.1109/TIT.1976.1055607}{doi:10.1109/TIT.1976.1055607}.
\bibitem{charpenay}
N.~Charpenay, M.~Le Treust, A.~Roumy. Side Information Design in Zero-Error Coding for Computing.
\emph{Entropy}, 26(4), 338, 2024.
\href{https://doi.org/10.3390/e26040338}{doi:10.3390/e26040338}.
\end{thebibliography}
}

\end{document}
